
% define some styles to underline the sematic of code fragments:
\newcommand{\obj}[1]{{\it #1}}		% identifier of a created Tcl-object
\newcommand{\cmd}[1]{{\tt #1}}		% a command of an object
\newcommand{\para}[1]{$<$#1$>$}		% an argument of a command call
\newcommand{\refarg}[1]{{\tt #1}}		% the argument mentioned within a text
\newcommand{\refslot}[1]{'{\tt #1}'}	% a Tcl-slot of an object mentioned with a text


\chapter{The Tcl-interface}
cb\_talk - 
Wrapper class for libcb functions to Tcl.

\section{Description}
	Wrapper class, for the ConceptBase application interface following
	the classic "talk" protocol for ConceptBase. Wrappes the methods
 \begin{description}
  \item[tell] to store something into the knowledge base
  \item[ask] ask a query of the knowledge base
  \item[untell] to transfer frames from the current to the historical database
 \end{description}
	connect and disconnect to the Tcl-Level. 
\section{Example}
 \label{tcl_example}
	An example of Tcl code that shows how the class can be used:
 \begin{verbatim}
  cb_talk server
   server connect 4223 gauguin example
    server tell "Class Employee in SimpleClass end"
   server disconnect
  server delete 
 \end{verbatim}
\section{Administration of the object}
 \subsection{Construction} 
  \label{tcl_construction}
	\begin{method}{\cmd{cb\_talk} [\obj{cb\_server}]}
	  \begin{description}
	   \item[ cb\_server]
		The command ``cb\_talk'' creates a Tcl object.
		The name of the new object can be determined by this 
		optional argument. % \refarg{cb\_server}.
	   \item[ return]
		If the name is not given the interface will generate
		a name in the form ``cb\_talk\_\_''+i, 
		where i is an integer. 
	  \end{description}
	If you want to use this object you must connect it with a 
	 running CBserver (see \ref{tcl_connection}).
	\end{method}

 \subsection{Destruction}
	\begin{method}{\obj{cb\_server} \cmd{delete}} 
	This function destroys the object. 
	You should use this function at the end of your application.
	If the object is connected then this function will end the 
	 connection.
	\end{method}

\section{Administration of the connection}
 \subsection{Connection} 
  \label{tcl_connection}
	If you want to work with a database then you have to connect 
	 the cb\_talk object with a running server.
	This is the purpose of the function connect: \\
	\begin{method}{
		\obj{cb\_server} \cmd{connect} \para{port} \para{host} 
		\para{client} [\para{user}] 
		}
%	The parametes are as follows:
	  \begin{description}	
%	     \item[cb\_server] This is the identifier of the cb\_talk object, 
%			which was created before.
	     \item[ port] The port number determines the channel for 
		the ipc-protocol. This number should be in the range of 
		2000 - 9999.
		The concret number should be given as parameter to the
		CBserver by the argument '-p' while starting it.
	     \item[ host] 
		The host name describes on which machine the CBserver 
		is running.
		If the server runs on the same machine as the client 
		then you can take 'localhost' as argument.
	     \item[ client] 
		This name will be used to register your client application
		within the CBserver.
	     \item[ user] This argument is optional. 
		It should contain the name of the current user. 
		The default value will be searched in
		the environment variables under 'USER'.
	     \item[ return]
		The return value is "1" if the connection is successful,
		else "0" will be returned.
		The handling of errors will be explained 
		in (see \ref{manageErrors}).
	  \end{description}
%	 and the error message is available by:
%	 \begin{verbatim}
%		cb_server GetErrorMsg
%	 \end{verbatim}
	\end{method}

 \subsection{Disconnection}
	At the end of a session with a special server it is useful to
	 disconnect the instance from the CBserver.
	If you delete the object it will be done automatically,
	 but if you want to connect it with an other server 
	 then you should use this function.
	The Tcl call is quiet simple: \\
	\begin{method}{\obj{cb\_server} \cmd{disconnect}} 
	  \begin{description}
	   \item[ return] There are no arguments necessary and 
		the return value is "1" if the disconnection was successful
	 	else "0" will be returned (see \ref{manageErrors}).
	  \end{description}
%	 and the error message is available by:
% 	 \begin{verbatim}
%		cb_server GetErrorMsg
%	 \end{verbatim}
	Notice that the instance is still existing,
	 though disconnecting it, so in normal cases this
	 command should be followed by: 
	 \begin{verbatim}
		cb_server delete
	 \end{verbatim}
	\end{method}

\section{IO-Access}
	Implements the {\tt cb\_talk} protocol to a ConceptBase server.
	The Tcl commands of this IO category will return a boolean value,
  	 ``1'' if command was successful, ``0'' if an error occurs.
  	In the last case it is usefull to manage the error, 
	 what will be described in \ref{manageErrors}. \\
	There are three fundamental functions on which the access is based,
	namely {\tt tell}, {\tt untell} and {\tt ask}.
	This functions should be fit to store some data into the 
	 knowledge base, delete it or retrieve any information from it. \\
	The precondition of this functions is that the object is connected
	 with a running server.
 \subsection{Tell}
	To store some information into the database of the server you
	 should use the function {\tt tell}. 
	These information can consists of any Telos frames such as
	 class definitions, individuals or queryclasses. 
	The Tcl call of this function has the following form: \\
	\begin{method}{\obj{cb\_server} \cmd{tell} \para{telos}} 
	  \begin{description}
	    \item[ telos] Must be a correct Telos frame.
%	The argument \refarg{telos} must be a correct Telos frame.
	    \item[ return] The return value 
		is "1" or "0" (see \ref{manageErrors}).
	  \end{description}
	\end{method}
 \subsection{Untell}
	The opposite of the {\tt tell} command is the {\tt untell} function.
	It deletes the information in the current database. \\
	\begin{method}{\obj{cb\_server} \cmd{untell} \para{telos}} 
	  \begin{description}
	    \item[ telos] Must be
	% The argument \refarg{telos} must be 
		correct Telos and the deleting of the object in 
		the knowledge base does not hurt any integration-constraint.
		In this case the will be moved from the actual database 
		into the historical knowledge. 
	    \item[ return] 
		If the delete transaction was successful 
		the return value is "1" else "0" (see \ref{manageErrors}).
	  \end{description}
	\end{method}
 \subsection{Ask}
	To retrieve any information you should use the {\tt ask } function: \\
	\begin{method}{\obj{cb\_server} \cmd{ask} \para{answerVar} 
			\para{query} [\para{format}] 
		}
%	This Tcl call asks the CBserver a \refarg{query} 
%	 and stores it to the answer.
	  \begin{description}
	    \item[ answerVar]
    		% The argument \refarg{answerVar} 
		This argument is the Tcl name of the variable 
	 	for the answer.
		%\\ Note that the value of the Tcl variable is overwritten 
		%in any case. Because if the Tcl variable does not exists, 
		%it should be ensured that it would be created, 
		%though an error occurs. \\
	    \item[ query] Must be a correct telos query.
		%The \refarg{query} must be a correct telos query.\\
	    \item[ format] 
		This argument is the optional format of
		% The \refarg{format} is the optional format of 
	 	FRAME, TELOS, FRAGMENT, LABEL or BOOL
			(see \ref{tcl_answerformats}).
		If the format is not given explicitly, 
		the method uses the default format in the 
		Tcl\_Slot \refslot{-answerformat} 
			(see \ref{tcl_slot_answerformat}).
	    \item[ return]
		The return value is "1" or "0" (see \ref{manageErrors})
	  \end{description}
	\end{method}
 \subsection{Formats} 
  \label{tcl_answerformats}
	These formats are for the answer of a query 
	 obtained by the ask command.
	The function lastAnswer gives the original text of the CBserver.
	But the answervariable of the ask command contains an object of the 
	wished format.
  \begin{description}	
   \item[FRAME] is the default format. \\
	The answer is the name of the Tcl object from the class cb\_ptelos.
	This object supports different functions to select an object 
	 in the frame and to retrieve any information about it.
    	Note that if the answer has a boolean value, e.g. 
    \begin{verbatim}
	<cb_server> ask <answerVar> "exists(\[ Integer/objname \])" FRAME
    \end{verbatim}
	The return value of the ask command is "0", because the answer is
	of the type BOOL
   \item[TELOS]
  	The answer is the original text in the TELOS format.
   \item[FRAGMENT]
  	The answer is the original text in the FRAGMENT format.
   \item[LABEL]
	The answer is a Tcl list of the object identifier of an answer frame.
   \item[BOOL]
	The answer is the "1" or "0" equivalent to the boolean 
	 representation of Tcl.
	If the answer was not a boolean value then the return value of
	 the ask command is "0".
  \end{description}
 \section{Last values}
	These methods supports a read-access on different variables 
	of the \obj{cb\_server} Tcl-object
  \subsection{Connection parameters}
%   \subsubsection{Port}
	\begin{method}{\obj{cb\_server} \cmd{lastPort}}
	  \begin{description}
	    \item[ return]
		Returns the value of the port of the last connection call.\\
		In the example \ref{tcl_example} this number 
		 would be ``4223''.
	  \end{description}
	\end{method}
%   \subsubsection{Host}
	\begin{method}{\obj{cb\_server} \cmd{lastHostname}}
	  \begin{description}
	    \item[ return]
		Returns the value of the host name 
		of the last connection call. \\
		In the example \ref{tcl_example} this name
		 would be ``gauguin''.
	  \end{description}
	\end{method}
%   \subsubsection{Client}
	\begin{method}{\obj{cb\_server} \cmd{lastClientname}}
	  \begin{description}
	    \item[ return]
		Returns the value of the client name of the 
		last connection call. \\
		In the example \ref{tcl_example} this value
		 would be ``example''.
	  \end{description}
	\end{method}
%   \subsubsection{User}
	\begin{method}{\obj{cb\_server} \cmd{lastUsername}}
	  \begin{description}
	    \item[ return]
		Returns the value of the user name of the
		 last connection call.
%		In the example \ref{tcl_example} this value
%		 corresponds to the environment value of ``USER''
	  \end{description}
	\end{method}
  \subsection{Access parameters}
%   \subsubsection{Answerformat}
	\begin{method}{\obj{cb\_server} \cmd{lastAnswerformat}}
	  \begin{description}
	    \item[ return]
		Returns the answer format of the last ask call. \\
		Note if no ask was executed, 
		 then the result is an empty string. \\
		Note that the last used format will be returned, 
		 independent if an answerformat was given as optional
		 argument of ask or as default value.
	  \end{description}
	 The purpose of this function is to validate, 
	 that a given answer has the right format.
	\end{method}
%   \subsubsection{Query}
	\begin{method}{\obj{cb\_server} \cmd{lastQuery}}
	  \begin{description}
	    \item[ return]
		Returns the query of the last ask call. 
	  \end{description}
	\end{method}
%   \subsubsection{Answer}
	\begin{method}{\obj{cb\_server} \cmd{lastAnswer}}
	  \begin{description}
	    \item[ return]
		Returns the original text of the answer
		 of the last ask call. 
	  \end{description}
		Note that the answer is the text,
		 which was returned from the CBserver.
	 	This means that a query with the answer format FRAME
		 stores internal a text in the TELOS format. 
		Or the answer format LABEL is internal
		 a list with elements separated by ``,''.
	\end{method}
%   \subsubsection{Answerobject}
	\begin{method}{\obj{cb\_server} \cmd{lastAnswerobject}}
	  \begin{description}
	    \item[ return]
		Returns the answer object which was stored 
		 to an answer variable by the last ask call. 
	  \end{description}
	Note if the answer was of the type FRAME, but the generated 
	 ``cb\_ptelos''-object was deleted, 
	 then this function will return the ident of this object,
	 but it has no access to it, because it doesn't exist anymore.
	\end{method}
  \subsection{Error} 
   \label{manageErrors}
	The following functions serve as error management
	The last completion value contains the error as short string constant.
	The error message contains the description of the occurred error in
	 human readable manner. \\
%   \subsubsection{Completion}
	\begin{method}{\obj{cb\_server} \cmd{lastCompletion}}
	  \begin{description}
	    \item[ return]
  		Returns the value of the completion status 
		 of the last operation. \\
		Note that the return value is a constant name equivalent 
		 to the names of ``CB\_interface.h'', like:	
		\begin{description}
		 \item[CB\_OK] The server access was successful
		 \item[CB\_ERROR] Anything was wrong, but it was non standard 
			so that the error message contains relevant information
		 \item[NOT\_HANDLED] The server doesn't know how to handle the
			occurred error
		 \item[CB\_TIMEOUT] Busywaiting for the server answer reaches
			the timeout limit
		 \item[CB\_CONN\_BROKEN] Signals that the connection was broken
		\end{description}
		Some cbtcl-interface specific constants are:
		\begin{description}
 		 \item[CB\_RECONNECTED] If the connection was broken, 
			but reconnected automatically. 
		 \item[CB\_UNKOWN\_COMPLETION] If the development 
			of ''CB\_interface.h'' has some new constants, 
			but it was forgotten to extend it
			in the module for the Tcl-interface.
		 \item[CB\_WRONG\_FORMAT] 
			If the answer was not of the expected type.
			For example a query was expected in the type BOOL, 
			 but it returns a list.
		 \item[CB\_NOT\_CONNECTED] 
			There was an attempt to gain access to the CBserver, 
			but there wasn't any connection.
		\end{description}
	  \end{description}
	\end{method}
%   \subsubsection{lastErrormessage}
	\begin{method}{\obj{cb\_server} \cmd{lastErrormessage}}
	  \begin{description}
	    \item[ return]
		Returns the error message of the last failure 
		 without deleting it. 
	  \end{description}
	The purpose of this function is to read the error message, 
	 without any side effect, in contrast to getErrorMsg, 
	 which deletes the error message.
	If a second error occurs the errormessage will be overwritten anyway.
	So it is not necessary to call \cmd{getErrorMsg},
  	 only if you want to mark that the error is handled by the program.
	\end{method}
  \subsection{Tcl-slots}
	The Tcl-slots are internal variables of the object.
	They can be determined optionally during creation of the object.
	For example: \\
	\cmd{cb\_talk} \obj{cb\_server} [\cmd{-answerformat} LABEL] 
					[\cmd{-rollback} Now] 
	\\
	Note that the number of configuration parameters are free.\\
	Also the object can be configured later, such as:\\
	\obj{cb\_server} \cmd{configure} \cmd{-answerformat} BOOL
	\\
	The 'configure' command behaves the same as the 'configure' command of
	 other Tcl objects.	
   \subsubsection{-errorhandler}
	These slot contains as default an empty script.
	You can store into this slot the name of a Tcl procedure 
	 with two argument. The first argument is the name of the server object
	 and the second is the completion value of the error.
	These procedure will be invoked by an error during the access to the
	 CBserver. 
	 The purpose of the function should be handle the occurred error.
	 The argument given to this error function contains the value
	  of {\tt lastCompletion}.
	 This is necessary because the completion value itself will be
	  temporarily set to ``CB\_OK'', so that you can work with the server
	  to analyze the error.
	 If the function returns ``1'' to signal that it handles the error
	 then the completion value will be set to ``CB\_OK'' else the original
	 completion value will be restored.
	A simple example for such an error handler:
  \begin{verbatim}
   proc CBaccessError { server orgCompletion } {
        if {($orgCompletion == "CB_WRONG_FORMAT") 
            && ([$server lastAnswerformat] == "BOOL")} then {
                # overwrites the answer as FALSE
                $server resetAnswer "0"
                $server resetAnswerobject "0"
                # mark the error as handled
                return "1"
        } else {
                # the error was not handled
                return "0"
        }	
   }
  \end{verbatim}
	
   \subsubsection{-rollback}
	The default rollback time for the access to the CBserver is ``Now''.
	 This exposes the current state of the database.
	But if you want query the state of the database of a historical date
	 then you can configure this value for example to ``day(1995,12,15)''.
   \subsubsection{-answerformat}
    \label{tcl_slot_answerformat}
	The default answer format is ``FRAME''.
	This format is a special object to handle a list of telos frames
	 as data structure, 
	 which will be explained later in section (\ref{cb_ptelos}).
	The other formats that are available are 
		FRAME, TELOS, FRAGMENT, LABEL or BOOL.
	They were mentioned in section (\ref{tcl_answerformats}).

	
 \section{Format FRAME} \label{cb_ptelos}
	If the answer format ``FRAME'' is used then the cb\_talk object
	generates an object of the type cb\_ptelos.
	The answer contains the name of these object, 
	 which will be handled as global Tcl variable.
	The generated name of the cb\_ptelos object will be of the form:
	 ``cb\_ptelos\_\_''+i, where i is an integer. \\
	
  \subsection{Description of cb\_ptelos}
	Wrapper class, for the ConceptBase application interface. 
	It allows parsing of telos strings and the access to the 
	the substructures of the frames, like classes, attributes, labels. 
%	\\
%	The syntax and semantic of the methods follows the C++ wrapper
%	{\tt parsedTelosObjects}, while the implementation is based
%	on the telos library xi/telos of Claudia Welter.

 \subsection{Example} 
  \label{ptelos_example}
	An example of Tcl code that shows how the class can be used:
  \begin{verbatim}
    # the following both lines will be done by the cb_server object implicit:
    cb_ptelos ptelos
    ptelos source "Class Employee in SimpleClass end
                   Class Manager isA Employee end"
    # the rest is to implement by application:
    set oid [ptelos GetObjectId]
    # Note: it is important to dispose the object at the end of the work:
    ptelos delete
  \end{verbatim}
 \subsection{Administration}
	The first both Tcl commands namely the constructor {\tt cb\_ptelos}
	and the method {\tt source} are used by the ask method of the 
	cb\_server object with answerformat FRAME. 
	So it is not necessary to call this functions directly. \\
	But the third Tcl command namely the destructor {\tt delete}
	 should be called when the work is done to free the memory. \\
% \subsection{Construction}
	\begin{method}{\cmd{cb\_ptelos} \obj{ptelos}}
	  \begin{description}
	     \item[ ptelos]
		Determines the name of the object optionally.
	     \item[ return] 
		Returns the name given to the object.
	  \end{description}
	The standard constructor initializes an empty instance. 
	This constructor is implicit used by the Tcl command:
	 \obj{cb\_server} \cmd{ask} .... ``FRAME''\\
	So in normal case you do not need to use this command.
	\end{method}

% \subsection{Parsing}
	\begin{method}{\obj{cb\_server} \cmd{source} \para{telos}}
	  \begin{description}
	     \item[ telos]
		Has to be a valid Telos string, 
		which should pe parsed.
	     \item[ return]
		In the normal case the result is ``1'', 
		but if the parser fails it will be ``0''.
	  \end{description}
%	The source method parses the argument. The argument has to be a
%	valid Telos string. 
	After the parsing you can gain access to
	the internal structure of the Telos string.\\
	In normal cases it is not necessary to use this function directly.
	\end{method}

% \subsection{Destruction}
	\begin{method}{\obj{ptelos} \cmd{delete}}
	Deletes the instance of the cb\_ptelos-object and 
	 the fragments of the last parsed answer.
	\end{method}
 \subsection{Select}	
	After a query in the ``FRAME'' format, when the telos answer will 
	 be parsed automatically, you can select one of the frames. 
	As a start the first frame is selected. \\

%  \begin{itemize}
	\begin{method}{\obj{cb\_ptelos} \cmd{first}} 
    	  \begin{description}
	     \item[ return] ``1'' or ``0'' as completion value.
	  \end{description}
	Selects the first frame in the parsed telos object.
	\end{method}

   \begin{method}{\obj{cb\_ptelos} \cmd{next}}
    	  \begin{description}
	     \item[ return] ``1'' or ``0'' as completion value.
	  \end{description}
	Selects the next frame in the parsed telos object.
	\end{method}

   \begin{method}{\obj{cb\_ptelos} \cmd{previous}}
    	  \begin{description}
	     \item[ return] ``1'' or ``0'' as completion value.
	  \end{description}
	Selects the previous frame in the parsed telos object.
	\end{method}

	\begin{method}{\obj{cb\_ptelos} \cmd{last}}
    	  \begin{description}
	     \item[ return] ``1'' or ``0'' as completion value.
	  \end{description}
	Selects the last frame in the parsed telos object.
	\end{method}
%  \end{itemize}
	The return values of the above mentioned functions is  normaly ``1'', 
	but if there are no frames on the new position it will be ``0'' 
	and the the current index does not change.

 \subsection{Status}
	There are two important functions to get the information of the
	 list dimension. The first returns the length of the list of frames.
	 The second returns the current index of the selected frame. \\
%  \begin{itemize}
	\begin{method}{\obj{cb\_ptelos} \cmd{count}}
    	  \begin{description}
	     \item[ return] 
		Returns the number of frames in the parsed telos object.
	  \end{description}
	  In the example \ref{ptelos_example} it would deliver ``2''.
	\end{method}
	\begin{method}{\obj{cb\_ptelos} \cmd{index}}
    	  \begin{description}	
	     \item[ return] 
		Returns the selected position of the current frame 
		 in the parsed telos object.
	  \end{description}
	\end{method}	
%  \end{itemize}
	Note that the index starts with 0 like other lists in Tcl.

 \subsection{Access}
	If you have selected a frame you can access the internals of it. 
	To access the attributes you have to specify a list of attribute 
	 categories. As a result you get a list of attribute labels or values 
	of the attributes that are in all the specified categories.
	\\

%	Note that the names chosen for the Tcl methods are the same
%	 as in the ``pTelosO.ch'' C++-header of Stefan Eherer.
%  \begin{itemize}
	\begin{method}{\obj{cb\_ptelos} \cmd{GetObjectId}}
    	  \begin{description}	
	     \item[ return] 
		Returns the object-identifier of the selected object 
		 if it exists, otherwise an empty string.
	  \end{description}
	In the example \ref{ptelos_example} the result is ``Employee''.
	If you call the next method before it would return ``Manager''.
	\end{method}	
   \begin{method}{\obj{cb\_ptelos} \cmd{Get\_IN\_OMEGA\_Class}}
    	  \begin{description}	
	     \item[ return] 
		Returns a Tcl string, which contains the omega-class 
		 of the current frame.
	  \end{description}
	In the example \ref{ptelos_example} the result would be ``Class''.
%	Other possibilities are ``Individual'' and ``Attribute''.
	\end{method}	
   \begin{method}{\obj{cb\_ptelos} \cmd{Get\_IN\_Classes}}
    	  \begin{description}	
	     \item[ return] 
		Returns a Tcl list, which contains the meta-class
		 of the current frame.
	  \end{description}
	In the example \ref{ptelos_example} the result 
	 would be ``SimpleClass''.
	\end{method}	
   \begin{method}{\obj{cb\_ptelos} \cmd{Get\_ISA\_Classes}}
    	  \begin{description}	
	     \item[ return] 
		Returns a Tcl list, which contains the super-class 
		 of the current frame.
	  \end{description}
	If you choice the second frame (``Manager'') 
	 of the example \ref{ptelos_example} by the next method 
	 then this function will deliver ``Employee''.
	\end{method}	
   \begin{method}{{cb\_ptelos} \cmd{Get\_Categories}}
    	  \begin{description}	
	     \item[ return] 
		Returns a Tcl list, which contains the categories
		 of the current frame.
	  \end{description}
	This function is useful to get a choice of possible filter categories
	 for the methods described in the next section.
%		(see \ref{ptelos_filtered_access}).
	\end{method}	
   \begin{method}{\obj{cb\_ptelos} \cmd{Get\_ValueOfLabel} \para{label}}
    	  \begin{description}	
	     \item[ label] 
		Is a string which should describe a label.
	     \item[ return] 
		Returns the value of the given \refarg{label} if it exists, 
		 otherwise the return is an empty string. 
	  \end{description}
	This enables a direct access on a concret value 
	 of the selected frame.
	\end{method}	
%  \end{itemize}

 \subsubsection{Get label and values filtered by categories}
%% \subsection{Filtered access by categories}
  \label{ptelos_filtered_access}
	The following three functions creating list of labels or values of
	 the selected frame. 
	The argument \refarg{searchedCategories} is a Tcl list of categories
	 that filters the labels and values, so that each label or value of
	 the resulting list is in all categories in the list.
	If the argument is an empty string then every label or value will be
	 listed in the order of appearence.	
	\\	
%  \begin{itemize}
   \begin{method}{\obj{cb\_ptelos} \cmd{GetLabelsOfCategories} \para{searchedCategories}}
    	  \begin{description}	
	     \item[ searchedCategories]
		Is a Tcl list of categories to filter the access 
		 as described above.
	     \item[ return]       
		Returns a list of the accordant labels.
	  \end{description}
	\end{method}	
   \begin{method}{\obj{cb\_ptelos} \cmd{GetValuesOfCategories} \para{searchedCategories}}
    	  \begin{description}	
	     \item[ searchedCategories]
		Is a Tcl list of categories to filter the access 
		 as described above.
	     \item[ return]       
		Returns a list of the adequate values.
	  \end{description}
	\end{method}	
   \begin{method}{\obj{cb\_ptelos} \cmd{GetLabValOfCategories} \para{searchedCategories}}
    	  \begin{description}	
	     \item[ searchedCategories]
		Is a Tcl list of categories to filter the access 
		 as described above.
	     \item[ return]       
		Returns a list of the correlating pairs of labels and values.
	  \end{description}
	Note these list is a Tcl list with sublists in the form
	 ``\{{\tt lab1 val1}\} \{{\tt lab2 val2}\} ...'', but it can easily 
	 decomposed by the standard Tcl commands for list such as \cmd{lindex}
		(see manpage for Tcl under the key ''list'').
	\end{method}	
%  \end{itemize}



% \section{AUTHOR(S)}
% Ralf Stoessel
% \section{COPYRIGHT}
% Copyright 1995 Lehrstuhl Informatik V - RWTH Aachen
